\textbf{Какие есть элементы, сопряжённые к а) $\sqrt{2}$; б)$\sqrt[3]{2}$; в) $\sqrt{2} + \sqrt{3}$; г) $\xi_p$ для простого $p$ над $\mathbb{Q}$?}

Воспользуемся тем, что сопряженный элемент должен быть корнем минимального многочлена для самого элемента. Каждый элемент сопряжен сам себе.

\textbf{а)} Мин. многочлен $x^2 - 2$. Один корень -- наш элемент $\sqrt{2}$, а второй -- ему сопряженный $-\sqrt{2}$.

\textbf{б)} Мин. многочлен $x^3 - 2$. Сопряженные исходному элементу корни $\omega\sqrt[3]{2}, \omega^2\sqrt[3]{2}$.

\textbf{в)} Многочлен $x^4 - 10x^2 + 1 = (x - \sqrt{2} - \sqrt{3})(x + \sqrt{2} - \sqrt{3})(x - \sqrt{2} + 
\sqrt{3})(x + \sqrt{2} + \sqrt{3})$. Несложно показать, что он неприводим над $\mathbb{Q}$ (а именно, подстановкой $y = x^2$ можно понять, что у него нет линейных делителей, а дальше -- от противного, анализируя коэффициенты у двух квадратных множителей можно получить противоречие). 

Этот многочлен -- минимальный для $\sqrt{2} + \sqrt{3}$. Соряженные исходному элементу корни $\sqrt{2} - \sqrt{3}, -\sqrt{2} + \sqrt{3}, - \sqrt{2} - \sqrt{3}$

\textbf{г)} Минимальный многочлен $x^{p-1} + \ldots + 1$ (из пункта 3-47). Его корни -- $\xi_n^1, \ldots, \xi_n^{p-1}$. Они и будут сопряженными элементами для $\xi_n$